\section{Problem Formulation}

The policy is any function $f_\theta$ that maps an observation to an action. It
must also report a per-step loss at each state-action pair it visits,
%
\begin{equation}
\label{eq:signal}
\ell^{(i)}_t \;=\; \mathcal{L}\big(f_\theta,\; s^{(i)}_t,\; a^{(i)}_t\big),
\end{equation}
%
where $i$ indexes an episode and $a^{(i)}_t$ is the action the policy itself
executed. The requirement stops there. For a diffusion policy $\mathcal{L}$ is the
denoising loss \cite{chi2023diffusionpolicy}; for a discrete policy it is the
negative log-likelihood of the action.

The policy is fitted to an initial set $\mathcal{D}_0$ by behaviour cloning
\cite{pomerleau1988alvinn,bain1995cloning}. Round $r$ opens by running the policy
on a fresh pool drawn from the task's reset distribution. The episodes it fails
form the round's failure set,
%
\begin{equation}
\label{eq:failures}
\mathcal{F}_r = \big\{\, f_i = (\tau_i,\; \ell^{(i)}_{1:T_i}) \;:\; \tau_i
\text{ is a failed episode of } f_\theta \,\big\}.
\end{equation}
%
The count $N = |\mathcal{F}_r|$ varies by round, and $T_i$ is the length of
episode $\tau_i$. The round closes by acquiring $D$ demonstrations from the expert
$\pi^\star$ and adding them to the training set \cite{ross2011dagger},
%
\begin{equation}
\label{eq:loop}
\begin{aligned}
\mathcal{D}_r &= \mathcal{D}_{r-1} \cup \{d_{r,1},\dots,d_{r,D}\},\\
\theta &\leftarrow \arg\min_\theta\ \mathbb{E}_{(s,a)\sim\mathcal{D}_r}
\big[\mathcal{L}_{\mathrm{BC}}(f_\theta, s, a)\big] \ \text{ for } r \equiv 0
\!\!\pmod m.
\end{aligned}
\end{equation}
%
The loop halts when $|\mathcal{D}_r| - |\mathcal{D}_0| = B$, so it runs $B/D$
rounds. Retraining runs at a per-task cadence $m$ set identically for every
method. Success is measured on a fixed set of test episodes, shared by every method and
seen by none while demonstrations are collected.

Every method here spends a round by the same three-part rule,
$A = (t^\star, C_{\mathrm{tgt}}, \xi)$. The step $t^\star$ is where the expert
takes over: the \emph{when}. The mode $C_{\mathrm{tgt}}$ is the failure the round
is spent on: the \emph{which}. The state $\xi$ is where the corrective
demonstration begins: the \emph{where}. A query gate computes $t^\star$ alone,
reading $C_{\mathrm{tgt}}$ and $\xi$ off whichever failed episode tripped its
threshold. The uniform-random control sets all three by sampling. DISEIL keeps
that same $t^\star$ and computes $C_{\mathrm{tgt}}$ and $\xi$ itself. Nothing else
in the loop changes, so the comparison is controlled.

\section{DISEIL}

\begin{figure*}[t]
\centering
\includegraphics[width=\textwidth]{figures/architecture.pdf}
\caption{DISEIL. Equation~\ref{eq:flag} returns the step $t^\star$ at which the
policy first becomes uncertain. A query gate computes that step and stops. DISEIL
continues from it. The perception branch meets the geometric branch only at the
prescription, so no model output enters the partition. ``Reasoning LLM'' and
``Prescription LLM'' are the two roles of the prescription model, a large language
model. ``KAG Grounding'' is the knowledge-augmented-generation constraint graph
$\mathcal{K}$. The supplementary specifies the solvability screen, drawn here as
the rollout and revise branch, where the policy is run once from the prescribed
start.}
\label{fig:architecture}
\end{figure*}

DISEIL runs four stages once per round (Figure~\ref{fig:architecture}). It
perceives the round's failures and partitions them into failure modes. It then
prioritises one mode and prescribes the $D$ demonstrations the expert is asked to
supply. Two models fill
three roles. The \emph{vision-language model} \cite{bai2025qwen3vl} reads rendered
frames. The \emph{prescription model} ($\mathrm{PM}$) \cite{yang2025qwen3} is a
text-only reasoning model that labels each failure's root cause and prescribes the
round's request. Neither model enters the partition.

\paragraph{Perceive.} A round opens by describing each failure twice, from the
same step of the same episode. The geometric description feeds the partition. The
language one feeds the prescription. The two never mix.
Each failure is anchored at one step: the first at which the loss crosses a
threshold $\eta$ and stays above it for $K$ consecutive steps,
%
\begin{equation}
\label{eq:flag}
t^\star_i = \min\big\{\, t : \ell^{(i)}_u > \eta \ \ \forall u \in [\,t,\,
t+K-1\,] \,\big\}.
\end{equation}
%
If the loss never crosses $\eta$, the anchor falls back to $\arg\max_t
\ell^{(i)}_t$. The threshold is a quantile of the training-loss distribution,
recalibrated at every retrain. Equation~\ref{eq:flag} is Diff-DAgger's rule, used
unchanged \cite{lee2025diffdagger}. The highest-loss step arrives late, so an
expert taking over there inherits a corrupted state with little episode left.

Each failure is then reduced to a six-dimensional vector, read from the
simulator's internal state at step $t^\star_i$. On the manipulation tasks it is
%
\begin{equation}
\label{eq:desc}
\begin{aligned}
\phi_i &= \big[\, p_{x,i},\; p_{y,i},\; \sin\psi_i,\; \cos\psi_i,\; \rho_i,\;
\delta_i \,\big],\\
\rho_i &= t^\star_i / T_i, \qquad
\delta_i = \big\| p^{\mathrm{ee}}_i - p^{\mathrm{obj}}_i \big\|_2.
\end{aligned}
\end{equation}
%
Here $(p_{x,i}, p_{y,i})$ is the planar position of the task-relevant object and
$\psi_i$ is its yaw. The yaw enters through its sine and cosine, so two yaws
either side of $\pm\pi$ are not falsely far apart. The fraction of the episode
completed before step $t^\star_i$ is $\rho_i$, and $\delta_i$ is the
end-effector-to-object distance. Each task refills the same six slots. The
descriptor is small because the failure sets are small: over a few dozen points,
weakly informative dimensions flatten the pairwise distances that decide which
clusters merge. The supplementary measures the effect.

In parallel, three rendered frames go to the vision-language model: one at the
episode's start, one at step $t^\star_i$, one at its end. The model returns a short
spatial account. The prescription model turns that account into a root cause and a
trajectory phase. Both come from a fixed vocabulary held in the constraint graph
$\mathcal{K}$ of the Prescribe stage below, not invented. The division of labour
follows the evidence: these models name a cause competently when handed structured
evidence \cite{duan2025aha,liu2023reflect}, but reason poorly about metric geometry
from pixels alone \cite{chen2024spatialvlm,fu2024blink}.

\paragraph{Partition.} A generic clustering step $\mathcal{C}$ groups the round's
failures, running over the standardised descriptors,
%
\begin{equation}
\label{eq:kstar}
\begin{aligned}
\tilde{X}_i &= \frac{\phi_i - \mu_\phi}{\sigma_\phi}, \qquad
\{C_1,\dots,C_{k^\star}\} = \mathcal{C}\big(\tilde{X},\, k^\star\big),\\
k^\star &= \arg\max_{k \in [2,\,k_{\max}]} \operatorname{sil}(k).
\end{aligned}
\end{equation}
%
Here $\mu_\phi$ and $\sigma_\phi$ are the per-dimension mean and standard deviation
of the round's own descriptors $\{\phi_i\}_{i=1}^{N}$, and the cluster count is
capped at $k_{\max} = \max(2, \min(6, N-1))$. The clustering step is generic by
design.
Agglomerative clustering is the choice made here \cite{ward1963hierarchical} and
k-means would serve \cite{lloyd1982kmeans}. The silhouette criterion picks
$k^\star$, used unmodified \cite{rousseeuw1987silhouette,pedregosa2011sklearn}.
DISEIL claims a partition step, not one implementation of it.

Each mode carries its centroid in raw pose coordinates; its mean peak loss
$\bar{L}_C$, averaged over its members' $\mathrm{peak}_i = \max_t \ell^{(i)}_t$;
and a representative $\mathrm{rep}(C)$, the member nearest the cluster mean in
standardised space. The dominant mode $C^\star$ holds the most members, ties
broken by mean peak loss. The peak scores a failure while the first crossing
locates it. Below four remaining failures each failure becomes its own
single-member mode, and Equation~\ref{eq:target} reduces to selecting the failure
of highest peak loss. A
mode is named by the majority root cause among its members, since reporting a
failure in mode 2 says nothing.

\paragraph{Prioritise.} Prioritise picks the target mode $C_{\mathrm{tgt}}$. The
Prescribe stage below picks the starting state $\xi$. Consider every mode whose
size is within one member of the largest.
The round goes to whichever of them has the highest mean peak loss,
%
\begin{equation}
\label{eq:target}
C_{\mathrm{tgt}} \;=\; \arg\max_{C \,:\, |C| \,\ge\, |C^\star| - 1} \
\bar{L}_C .
\end{equation}
%
The size constraint keeps the target inside the bulk of the round's failures, so a
mode that barely exists cannot capture the budget on one badly failed episode.
Because one demonstration rarely removes a mode outright, a rule of this shape can
return the same mode round after round. DISEIL therefore has a cluster memory,
which discounts modes already corrected. The memory is a configurable
task-specific component outside the core rule, on in every run here. The
prescription model is then shown at most $\kappa$ cited failures from the target
mode: the representative first, then the worst-loss failure, then farthest-point
selection \cite{eldar1997fps} to span the mode.

\paragraph{Prescribe.} The prescription model receives three inputs: the target
mode's anchor geometry, the cited failures with their root-cause labels, and the
task's constraints. It returns the round's request, with a confidence score and a
rationale. Each requested demonstration takes one of two forms. A \emph{targeted
correction} names one cited failure; that exact episode is re-instantiated and the
expert takes over at the anchored step. A \emph{bridging placement} names two or
three cited failures and asks for a new configuration between them, from which the
expert demonstrates a complete episode. Bridging lets a prescription be easier than
any failure it addresses. Suppose a mode lies far outside what the policy can
solve. A targeted correction there asks for a large jump; a bridged one asks for a
step the policy can absorb. The graded-placement argument
is reverse curriculum generation's \cite{florensa2017reversecurriculum}, applied to
a demonstration request rather than a training reset.

A model asked for a configuration of the world will sometimes name one the world
cannot produce. The constraint graph $\mathcal{K}$ makes such a request checkable.
It holds workspace bounds, spawn ranges, reachability, controller limits, the
success predicate and the failure-mode vocabulary as structured knowledge. It is a
robot knowledge base \cite{tenorth2013knowrob}, not a retrieval index over text
\cite{lewis2020rag,edge2024graphrag}. Verification is a loop. The model proposes a
configuration. A map $g$ turns it into a reset specification $\xi$. A test $V$
checks $\xi$. If $\xi$ fails, the violation returns as feedback and the model
proposes again,
%
\begin{equation}
\label{eq:verify}
\begin{aligned}
\mathrm{cmd}^{(j)} &= \mathrm{PM}\big(C_{\mathrm{tgt}},\ S,\ \mathcal{K},\
\text{violation}(\xi^{(j-1)})\big),\\[1pt]
\xi^{(j)} &= g\big(\mathrm{cmd}^{(j)}\big),\\[1pt]
V(\xi) &= \mathbf{1}\big[\xi \in \mathcal{W}_{\mathcal{K}}\big] \wedge
\mathbf{1}\big[\mathrm{reachable}_{\mathcal{K}}(\xi)\big]\\[1pt]
&\quad \wedge\ \mathbf{1}\big[\mathrm{valid\text{-}path}_{\mathcal{K}}(\xi)
\big].
\end{aligned}
\end{equation}
%
The workspace bounds are $\mathcal{W}_{\mathcal{K}}$, and the conjuncts of $V$ are
the constraints $\mathcal{K}$ holds for that task. There is no feedback on the
first attempt, so $\text{violation}(\xi^{(0)}) = \emptyset$. On the grid task the
last conjunct instead checks that a path exists, using A\textsuperscript{$\star$}
and breadth-first search \cite{hart1968astar,cormen2022algorithms}, neither of
which ever acts as the expert. A failed attempt consumes no budget: the budget
counts demonstrations collected, not prescriptions proposed. After $J_{\max}$
attempts the round falls back to the nearest untried recorded failure. A second
screen rejects a prescription the policy already solves. The supplementary
specifies it, and no run below exercises it.

Every stage bears on what the policy is trained on and none on how it is fitted,
so the evaluation varies the acquisition rule and holds the rest of the loop fixed.
The supplementary states the round as pseudocode.
